\section{Method}

Reverse Edit Reconstruction evaluates explanation faithfulness by asking whether a structured edit can be recovered from an explanation. The method has four stages.

\paragraph{Edit extraction.}
Given a source sentence and a model prediction, we extract token-level edits. The current implementation uses a deterministic token-diff fallback for executable scaffolding. Final experiments should use ERRANT or another verified GEC edit extractor to obtain linguistically meaningful boundaries and error types.

\paragraph{Explanation conditions.}
To avoid reducing reconstruction to answer copying, explanations should be evaluated under multiple conditions: explicit-edit explanations, implicit explanations, masked-target explanations, rule-only explanations, and minimal sufficient explanations. Additional controls include source-only, explanation-only, shuffled explanation, and raw edit string inputs.

\paragraph{Reconstruction.}
The reconstructor predicts whether the edit is reconstructable and, if so, outputs span boundaries, source text, target text, operation, and error type. A simple structured baseline can extract explicit phrases such as ``replace A with B'' from explanations. More capable reconstructors may use sequence-to-structure models or LLM-based structured extraction, but they must be compared under leakage controls.

\paragraph{Scoring.}
The reconstructed edit \(\tilde{e}\) is compared with the model-produced edit \(e\). We compute span exact match, span precision/recall/F1, source text match, target text match, operation accuracy, error type accuracy, and full edit exact match. For faithfulness classification, reconstruction agreement can be thresholded or used as features for a classifier. Negative examples are required to test whether the method rejects explanations with wrong location, wrong target, wrong direction, wrong type, swapped edit, or generic rationale.

